\documentclass[11pt,a4paper,headsepline,twoside,toc=bibliography]{scrreprt}

\usepackage[utf8]{inputenc}
\usepackage[T1]{fontenc}
\usepackage[latin]{armtex}
\usepackage[ngerman]{babel}
\usepackage{hyperref}
\usepackage{graphicx}
\usepackage{caption}
\usepackage{subcaption}
\usepackage{shorttoc}
\usepackage{booktabs}
\usepackage{listings}
\usepackage{todonotes}



% Styling
\pagestyle{headings}
\renewcommand{\chaptermark}[1]{\markboth{#1}{}}
\renewcommand{\sectionmark}[1]{\markright{#1}}

\usepackage{tipa}

\hypersetup{linktoc=all,hidelinks}

\newcommand{\HRule}{\rule{\linewidth}{0.5mm}}

\author{Mariam Hemmer}
\title{Interkulturelle Wahrnehmung von Emotionen in Sprache am Beispiel armenischer und deutscher Hörer}

\usepackage{csquotes}
\usepackage[
backend=biber,
style=apa,
alldates=year,
isbn=false,
url=false,
uniquename=init
]{biblatex}


\addbibresource{Masterarbeit.bib}

%\bibliography{Masterarbeit} 

\begin{document}
\begin{titlepage}
\pagenumbering{roman}


\begin{center}
{
\scshape

\includegraphics{siegel}\\[1cm]    

{\LARGE Ludwig-Maximilians-Universität München}

Fakultät für Sprach- und Literaturwissenschaften

Institut für Phonetik und Sprachverarbeitung\\[0.75cm]

{
\HRule \\[0.4cm]

{
\huge \bfseries Interkulturelle Wahrnehmung von Emotionen in Sprache am Beispiel armenischer und deutscher Hörer\\[0.3cm] \large Der Einfluss der Grundfrequenz und des Sprachtempos auf die Wahrnehmung der Emotionen\\[0.4cm]
}

\HRule \\[0.75cm]
}

{
\Large
Masterarbeit\\
Betreuer: PD Dr. Christoph Draxler\\[0.5cm]
Vorgelegt am 20.07.2020\\von Mariam Hemmer\\[0.5cm]
}

Matrikelnummer: 11168853

Hauptfach: Phonetik und Sprachverarbeitung

Adresse: Rembrandtstraße 6, 81245 München

Geburtsdatum: 20.05.1988

}
\end{center}
\cleardoublepage
\end{titlepage}

\chapter*{Erklärung der Eigenständigkeit}

Hiermit erkläre ich, dass ich die vorliegende Arbeit selbstständig und ohne Benutzung anderer als der angegebenen Hilfsmittel angefertigt habe. Alle Stellen, die wörtlich oder sinngemäß aus anderen veröffentlichten und nicht veröffentlichten Quellen entnommen wurden, wurden als solche kenntlich gemacht. Die Arbeit ist in gleicher oder ähnlicher Form oder auszugsweise im Rahmen einer anderen Prüfung noch nicht vorgelegt worden.

\vphantom{M}

\noindent München, den 20.07.2020

\vphantom{M}

\vphantom{M}

\begin{flushright}
$\overline{~~~~~~~~~~~~\mbox{\strut Mariam Hemmer}~~~~~~~~~~~~}$
\end{flushright}


\chapter*{Zusammenfassung}

\chapter*{Abstract}


\cleardoublepage

%\shorttoc{Overview}{0}

\tableofcontents
\cleardoublepage

\pagenumbering{arabic}


\chapter{Notizen}
\label{sec:general_notes}
\section{Analysen emotionaler Prosodie, Beate Wendt, 2007}
\begin{itemize}
	
		
	\item Seite 58:
	\begin{itemize}
		\item Es gibt eine muttersprachliche Kompetenz komplexe prosodische Muster für emotionsalisierte Ausdrücke abzurufen und auszuführen (Neuber 2002)
		\todo{Neuber 2002 missing}	
		\item überindividuelle Beziehungen zwischen Form und Funktion -> bestimmte Formmuster innerhalb einer Sprachgemeinschaft werden zu bestimmten Bedeutungkomponenten zugeordnet (Neuber 2002)	
	\end{itemize}

	\item Seite 59:
	\begin{itemize}	
		\item Hörexperiment/Wahrnehmungsexperiment - bottom-up vs. top-down Wahrnehmungsprozess 
		\item Analytisches Hören (Viregge 1996) ist relativ (Bose 2003)
		\todo{Viregge 1996, Bose 2003 missing}
	\end{itemize}

	\item Seite 64:
	\begin{itemize}
		\item Die Sprachgeschwindigkeit - keine wichtige Information
	\end{itemize}
	 
	
	\item Seite 70:
	\begin{itemize}
		\item Zur Sprachgeschwindigkeit:\\
		Laut (GWdA 1982) Koartikulation->Leitungsabfall wird in Atem audgeglichen, Leistungssteigung in Artikulation, das ergibt Assimilation und Dessimilation. \todo{ME: Leitungsabfall und Leistungssteigung sind sehr verdächtige Formulierung, lieber nicht benutzen!}
		\todo{Was ist GWdA 1982? Missing anyway!}	
	\end{itemize} 
\end{itemize}


\chapter{Einführung}
\label{sec:einfuehrung}
Menschen reagieren ständig. Die Reaktionen der Menschen sind oft mit Emotionen verbunden. Entweder werden Emotionen durch Reaktionen ausgelöst, oder die Reaktionen entstehen durch ausgelöste Emotionen: die Kontexte sind verschieden, sowie die Theorien (\todo{die Autoren in Psycholinguistik/Neurophysiologie Folien!})
Haben kulturelle Faktoren Einfluss auf die Emotionswahrnehmung? Humor in einer Fremdsprache wird nicht immer las humorvoll wahrgenommen. Zugrunde dessen können verschieden Faktoren liegen. So können kulturell eingeprägte moralische Vorstellungen dararuf einen Einfluss haben, ob ein Witz (\cite{Nevo2001}). Die Studie zeigt, dass Singapurische Studenten konservative Einstellung bezüglich des Inhalts der Witze gezeigt haben. Dahingegen haben amerikanischen Teilnehmern eher Witze gefallen, die agressive Inhalt beinhalteten. In einer verbalen Kommunikation können bestimmte akustische Reize eine Emotion auslösen. So wird Lachen als ansteckend bezeichnet (\cite{Provine1992}). Oder eine Geschichte, die als traurig eingeschätzt wird, kann die Emotion Trauer auslösen. Menschliche verbale Kommunikation kann also sowohl auf der semantischen als auch suprasemantischen Ebene eine Emotionsauslösende Wirkung haben.\\



Bei dieser Arbeit geht es um interkulturelle Wahrnehmung von Emotionen am Beispiel von deutschen und armenischen Hörern. Dabei wird folgendes untersucht:
\begin{itemize}
	\item die Übereinstimmung der Wahrnehmung von Emotionen in der verbalen Kommunikation bei den deustchen und armenischen Hörern 
	\item der Einfluss der akustischen Merkmale Grundfrequenz und Tempo auf die Emotionswahrnehmung 
\end{itemize}


\chapter{Stand der Forschung}
\label{sec:theory}

Interkulturelle Emotionswahrnehmung wird in vielen Kontexten durchgefürt. 

 \todo{don't forget to cite!}
\begin{itemize}
	\item Gesichtsausdruck
	\item Semantik - schriftliche emotionale Ausdrücke
	\item verbale bedeutungsvolle ausdrücke
	\item nicht verbale ausdrücke  
	 
\end{itemize}


Studien zeigen, dass (\cite{Pell2020})


\section{Deutsch - Armenisch: eine sprachliche Gegenüberstellung}
\label{sec:language_description}
Die moderne armenische Sprach besteht aus Ostarmenisch und Westarmenisch. Armeniens Amtssprache ist Ostarmenisch. Im Rahmen dieser Arbeit wird 


\section{Interkulturelle Zusammenhänge in der Emotionswahrnehmung}
\label{sec:crosscultural_emotion}

In diesem Abschnitt werden die Ansätze in der Forschung hinsichtlich der interkulturellen Emotionswahrnehmung dargestellt.Länder unterscheiden sich unter Anderem aufgrund ihrer Kultur. Die Sprache als Teil der Kultur spielt hier eine sehr große Rolle. Es gibt viele Definitionen für "Kultur". Definition!!! Eine weitverbreitete Definition ist: .... . Hofstede's defition. Auch Ansätze der Forscher, die interkulturelle Studien durchführen, unterscheiden sich. Edward T. Hall (1976) zum Beispiel unterscheidet die Kulturen abhängig von der Art der Kommunikation. Dahingegen werden von manchen Autoren die Kulturunterschiede auf einer mehrdimensionalen Ebene untersucht (\cite{Matsumoto}). Die Dimensionen, die eine Kultur beschreiben sollen, wurden von Hofstede (1980, 1983) vorgeschlagen, da Kulturunterschiede auch innerhalb eines Landes existieren können. Hofstedes Definitionen auflisten und beshcreiben. Hofstede (2001) definiert die Kultur als \emph{"the collective programming of the mind, that distinguishes the members one group or category of people from another."} (\todo{cite Hofstede 2001, s. 9}). Dazu Kritik: Baskeville - keine Kultur, sondern nur Wirtschaftliche Faktoren. 


In vielen Studien werden interkulturelle Vergleiche zwischen Ländern durchgeführt. Der Ansatz, dass Land und Kultur gleichgesetzt werden, kann in vielen bisherigen Studien gesichtet werden\todo{Warum werden Land und Kultur gleichgesetzt?} (Ekmann, usw.). Manche Experimente haben gezeigt, dass die Wahrnehmung der Emotionen aufgrund von Gesichtsausdrücken vorwiegend kulturübergreifend ist (Ekman 1972, Ekman and Friesen 1971, Ekman, Friesen and Ellsworth 1972, Ekman, Sorenson and Friesen 1969, Izard 1971). In einer weiteren Studie hat Ekman (1987) beschrieben, dass Teilnehmer verschiedener Kulturen die Intensität der Emotionen unterschiedlich bewertet haben. Eine andere Studie konnte feststellen, dass die Intensitätswahrnehmung der Emotionen von den linguistischen Bedeutungsunterschieden in den Kulturen nicht abhängig ist (Ekman and Matsumo in Press?). 


\citeauthor{Burkhardt2006} hat die interkulturelle Wahrnehmung von Emotionen  in Frankreich, Deutschland, Griechenland und in der Türkei untersucht \parencite{Burkhardt2006}. Dabei wurden den Teilnehmern semantisch bedeutungsvolle Sätze  präsentiert. Durch die systematische Manipulation der akustischen Merkmale Grundfrequenz, Jitter und Dauer sollte der Einfluss dieser Merkmale auf die Emotionswahrnehmung in oben genannten Ländern untersucht werden. Für jede Sprache konnten merkmalsabhängige Variationen bei der Diskrimination der Emotionen beobachtet werden, nichtsdestotrotz wurden keine signifikanten sprachabhängigen Unterschiede in der Emotionswahrnehmung festgestellt. Dies wird unter anderem dadurch begründet, dass die semantisch bedeutungsvollen Sätze nicht identisch sind, da sie für jedes Land in die jeweilige Sprache übersetzt werden mussten.




Die Beseitigung der linguistischen Faktoren kann in der Forschung dadurch gelöst werden, dass die Semantik in einem verbalen Ausdruck möglichst minimiert wird. Als solches werden oft nonsense-speech, Pseudowörter sowie non-verbale Ausdrücke, wie Husten, Lachen oder Weinen,  benutzt. Dabei wird es davon ausgegangen, dass die Emotionen durch bestimmte Musterprofile in die gesprochenen Sprache encodiert werden. Um herauszufinden, ob solche encodierte Emotionen in unabhängig von Sprache und Kultur erkannt werden können, hat \citeauthor{Scherer2001} ein Perzeptionsexperment in 9 unterschiedlichen Ländern durchgeführt. Dabei wurden Lautkombinationen aus sechs europäischen Sprachen (German, English, French, Italian, Spanish, and Danish) zusammengesetzt. Daraus wurden Sätze formuliert, die von Schauspielern und Schauspielerinnen in siben prosodischen Ausdrucksweisen ausgesprochen wurden. Neben neutral wurden die Basisemotionen Wut, Freude, Ekel, Angst, Trauer sowei  benutzt. \citeauthor{Scherer2001} weist darauf hin, dass es zwar kleine kultuabhängige Unterschiede festgestellt wurden, die Accuracy der Emotionserkennung in allen Ländern höher war, als es durch Zufall möglich wäre.   

Bei einer anderen Studie wurden die Interpretationen der emotionalen Ausdrücke untersucht (\cite). Wie die Studie abgelaufen ist und was dabei rausgekommen ist.


Hörer können die Emotionen besser zuordnen, wenn diese in eigener Sprache ausgesprochen werden (\cite{Pell2020, Chronaki2018}). 


Bei dieser Studie geht es um die Emotionserkennung in der gesprochenen Sprache. Aus diesem Grund sollte die gemeinsame Sprache einer Gruppe im Rahmen dieser Arbeit als Bedingung dafür dienen, diese Gruppe als ?? einer Kultur. Hinsichtlich der sprachlichen Aufteilung werden die Teilnehmer dieser Studie aufgrund ihrer Muttersprache zu einer Kultur zugeordnet.  

\section{Sprachperzeption}
\label{speech_perception}

Sprache ist das Hauptmittel der menschlichen Kommunikation. Unter Kommunikation versteht man Informationsaustausch. Dabei handelt es sich um einen Prozess, an dem mindestens zwei 'Parteien' beteiligt sind. Einer erzeugt die Information und gibt weiter an den Partner, der wiederum nimmt die Information wahr und verarbeitet diese. Bei der verbalen Kommunikation sind die Perteien dieses Prozesses der Sprecher und der Hörer, wobei der Sprecher gibt encodierte Informationen durch die verbalen Ausdrücke an den Hörer weiter, der zunächst diese Informationen perzepiert und anschließend dekodiert. Wenn unter Dekodierung der Informationen in einem verbalen Ausdruck dessen Zerlegung in Infromationsklassen verstanden wird, können diese Klassen wie folgt aussehen:

\begin{itemize}
	\item nicht-verbale: Informationen über den Sprecher, wie Alter, Geschlecht, Herkunft usw.
	\item semantische (oder eher verbale linguistische)
	\item parasemantische (oder eher verbale paralinguistische)
\end{itemize}

Das menschliche Gehör besteht aus drei Teilen: Außenohr, Mittelohr und Innenohr (hier kommt das bekannte Bild). Das Außenohr dient als Antenne, welche die empfagenen akustischen Reize verstärkt und an das Mittelohr überträgt. Hier werden die Reize durch Ambus und Hammer an das Innenohr übertragen. Die Verarbeitung der Frequenzen im Innenohr übernimmt die Cochlea - die Hörschnecke (hier kommt ebenfalls ein bild). Detaillierte Schilderung der Membrane und Bahnen und Haarzellen folgt noch! Unterschiedliche Teile der Cochlea sind für unterschiedliche Frequenzbereiche zuständig (hier kann natürlich auch ein bild der Cochlea mit der Zuteilung der Frequenzbereiche). Die durch die Basilarmembran empfangenen Frequenzen werden in ein elektrisches Signal umgewandelt und als Nervenimpuls and das Gehirn weitergeleitet. Hier können die Hörbahnen erklärt werden. Anschließend reagieren die Gehirnareale, die an der Sprachperzeption beteiligt sind, auf den Impuls und verarbeiten diesen. So wird zum Beispiel die Phonemerkennung im Brocca-Areal durchgeführt, während im Wernicke-Areal nach der Bedeutung dieser Lautkombination gesucht wird. \\

Es ist schon lange bekannt, dass die semantische Verarbeitung der Sprache weitgehend in der linken (oder eher in der dominanten) Hemisphäre durchgeführt wird. Dahingegen werden suprasemantische Ebenen der gesprochenen Sprache in der rechten oder nicht dominanten Gehirnhälfte verarbeitet. Zur Suprasemantik der gesprochenen Sprache zählt auch die Emotionswahrnehmung, die in dieser Arbeit untersucht wird. Die semantische Verarbeitung der Sprache kann einen Einfluss auf die Emotionswahrnehmung haben. So kann ein Wort mit einer positiven Bedeutung mit einer negativen Emotion ausgesprochen werden. Die Bedeutung des Wortes kann in diesem Fall als Störfaktor für die Emotionswahrnehmung bezeichnet werden, da die Semantik einen Einfluss auf die Verarbeitung der emotionalen Information hat, die durch Prosodie vermitelt wird. Für diese Arbeit steht die Emotionswahrnehmung im Fokus. Daher sollte die Semantik der Stimuli als Störfaktor betrachtet werden. Für diesen Zweck sollten also Sprachaufnahmen als Stimuli in Frage kommen, die keine Bedeutung haben. Da im Rahmen dieser Arbeit die interkulturelle Wahrnehmung der deutschen und armenischen Hörer untersucht wird, sollten die Stimuli sowohl auf Armenisch als auch auf Deutsch keine Bedeutung haben. Die Anforderungen dieser Studie für den Sprachkorpus können wie folgt zusammengefasst werden: 


\begin{itemize}
	\item der Sprachkorpus enthält emotionale Sprache
	\item mindestens ein Sprecher und eine Sprecherin
	\item der selbe Text ausgesprochen in allen Emotionen, die im Sprachkorpus enthalten sind
	\item mindestens folgende Basisemotionen : Freude, Wut, Trauer, Angst und neutral
	\item die Stimuli sollen möglichst kurz sein, eher Wörter als Sätze, da je länger der Stimulus, desto mehr  Informationen ist im Sprachsignal enhalten. Daher kann zu viel Information als Störeffekt für die Prosodieverarbeitung hinsichtlich der Emotionen sein.
\end{itemize}


Das Bayerische Archiv für Sprachsignale (BAS: https://www.bas.uni-muenchen.de/Bas/BasHomedeu.html) stellt Sprachdatenbanken zur Verfügung, die für die Forschung benutzt werden können. \\


Die akustischen Reize werden durch das Ohr wahrgenommen. Bereits das Außenohr verstärkt den empfangenen Schall und überträgt durch das Trommelfell (Ambus und Hammer) \\


\section{Emotion in Sprache}
\label{sec:emotion_in_speech}



\emph{Klassifizierung von Emotionen}
-Es gibt 3 Modelle der Klassifizierung der Emotionen\\

\emph{Basisemotionen}
Der Magdeburger Sprachkorpus beinhaltet die Basisemotionen Freude, Trauer, Wut, Ekel, Angst sowie neutral. Im Sprachkorpus sind die neutralen Sprachaufnahmen als 'sachlich' vermerkt (\cite{Wendt2002}). Im Sprachkorpus befanden sich deutsche Wörter, sie in diesen Basisemotionen von einem Schauspieler und von einer Schauspelering gesprochen wurden. Außerdem wurden aus diesen Wörtern Pseudowörter durch Umpositionierung der Laute erstellt, diese ebenfalls von den selben Schauspielern in obengenannten Emotionen ausgesprochen wurden. Die kreierten Pseudowörter entsprachen der deutschen Phonotaktik. In einem Perzeptionsexperiment wurden die ausgesprochenen Emotionen evaluiert (\cite{Wendt2003}). \\

Emotionen haben einen großen Einfluss auf die Grundfrequenz. Wenn der Sprecher traurig ist, ist seine Grundfrequenz eher iedriger, als wenn der selbe Sprecher glücklich oder wütend ist. Das hat physiologische Gründe. Zum Beispiel bei Wut oder Freude werden bestimmte Muskeln angespannt, die an der Sprachproduktion teilnehmen. Durch die Anspannung der Muskeln ändert sich die intrinsische Grundfrequenz. \\

-Komponenten-Prozess-Modell 
\todo[inline, color=red!40]{Scherer 1990, Paeschke 2003 missing!}
Reizbewertungsprüfung hat Einfluss auf das Nervensystem. Dadurch entsteht eine ganz bestimmte Veränderung der Muskelanspannung: diese Änderungen betreffen auch die Muskeln der Stimmproduktion (Paeschke 2003).\\

Dieser Vorgang betrifft die Fälle, wenn die Emotionen auf einer natürlichen Weise entstehen und einen Einfluss auf die Sprachproduktion haben. Jetzt ist die Frage, ob es möglich ist Emotionen zu spielen und diesen natürlichen Vorgang der Prosodieänderung bei der emotionalen Sprache nachzuahmen. 

\emph{ABER}: die Stimuli oder die Emotionen der Stimuli sind nicht durch eine natürlicheReizbewertungsprüfung der Sprecher in Stande gekommen, sondern wurden diese durch die Schauspieler entsprechend ihrer subjektiven Vorstellung produziert. Diese Stimuli wurde dennoch mit einem Perzeptionsexperiment evaluiert.\\



\section{Prosodie}
\label{sec:prosody}

Die menschliche Kommunikation ist vielschichtig. Die gesprochene Sprache enthält einerseits den konkreten Inhalt, was gesagt wurde, andererseits vermittelt sie ebenso die Einstellung und das Verhältnis des Sprechers zum Inhalt des Gesagten - also wie etwas gesagt wurde. Oft ist die Information darüber, wie etwas gesagt wurde, sogar viel wichtiger für den Hörer, als der Inhalt an sich (\textcite{Wildgruber2006}). Hinsichtlich der verbalen Kommunikation wird diese suprasemantische Information in der gesprochenen Sprache encodiert und durch akustische Cues vermittelt. Die Stimmqualität, der Intonationsverlauf der Stimme, sowie Variationen der Tonhöhe und der rythmischen Struktur, wie zum Beispiel Sprechtempo, geben Hinweise über den emotionalen Zustand des Sprechers. Diese Hinweise dienen für den Hörer als akustische Korrelate bei der Wahrnehmung der affektiven Prosodie \todo{zitieren: (Scherer 1986; Scherer et al. 2009; Sidtis \& Van Lancker Sidtis 2003}. 

Die Details hierzu werden im nächsten Kapitel behandelt. \\  

Bei der Prosodieverarbeitung im Gehirn ist die rechte Hemisphäre mehr beteiligt, als die linke Hemisphäre \parencite{Ackermann2014c} .  

\section{Akustische Merkmale}
\label{sec:prosody_features}


Welche Merkmale können die Emotionswahrnehmung beeinflussen? Wie werden diese Merkmale wahrgenommen und verarbeitet? Sprachperzeption ist ein komplexer Prozess, in dem mehrere Areale sowohl in der rechten als auch in der linken Hemisphäre des Gehirns beteiligt sind. Nichtsdestotrotz werden die wahrgenommenen Reize vorwiegend einseitig verarbeitet, da die Hemisphären nicht gleichermaßen für alle Verarbeitungsprozesse spezialisiert sind. Auch bei der Sprachverarbeitung findet eine Lateralisierung der Informationsverarbeitung im Gehirn statt. Zur Funktionalität der linken Hemisphäre zählen zum Beispiel die syntaktische und die phonologische Informationsverarbeitung, sowie andere sprachliche Fähigkeiten**. Im  Gegensatz dazu ist die rechte Hemisphäre zuständig für den Intonationsverlauf, ebenso wie für weitere suprasemantische Aspekte der Sprachverarbeitung. Der Intonationsverlauf kommt durch die Variation der Grundfrequenz zustande. Die Grundfrequenz, als ein sehr wichtiger akustischer Merkmal, vermittelt Informationen darüber, wer der Sprecher ist, und wie der Sprecher spricht. Außer Hinweisen über Geschlecht und Alter des Sprechers beinhaltet die Grundfrequenz auch Informationen über die Intentionen und den emotionalen Zustand des Sprechers. \\

Eine Studie zeigt, dass der Verarbeitungsprozess der Grundfrequenz durch die Lateralisierung der rechten Hemisphäre , in der auch die Emotionswahrnehmung stattfindet, zugeteilt wird (\cite{VanLancker1973a}). \\





Die Produktion und die Perzeption der emotionalen Sprache wird im Wesentlichen durch drei Merkmale beinflusst (\cite{Scherer2001})

\begin{enumerate}
	\item Grundfrequenz: Variation des Intonationsverlausf bei unterschiedlichen Emotionen 
	\item Sprachtempo: Bei manchen Emotionen redet man schneller oder langsamer, ein Beispiel wäre Traurigkeit, wo das Tempo verlangsamt wird, im Gegensatz zu wütend, wo man schneller redet
	\item Intensität
\end{enumerate}


\begin{figure}
	\centering
	\begin{subfigure}{1\textwidth}
		\centering
		\includegraphics[width=0.8\linewidth]{pics/sawet_s.png}
		\caption{neutral}
		\label{fig:sawet_s}
	\end{subfigure}
	\par\bigskip
	\begin{subfigure}{0.8\textwidth}
		\centering
		\includegraphics[width=1\linewidth]{pics/sawet_T.png}
		\caption{Trauer}
		\label{fig:sawet_T}
	\end{subfigure}
	\par\bigskip
	\begin{subfigure}{0.8\textwidth}
		\centering
		\includegraphics[width=1\linewidth]{pics/sawet_F.png}
		\caption{Freude}
		\label{fig:sawet_F}
	\end{subfigure}
	\caption{\textbf{Intonationsverlauf für verschiedene Emotionen.} das Wort \emph{sawet} (IPA: \textipa{[zavEt\super{h}]}) ausgesprochen von dem selben Sprecher (männliche Stimme) in prosodischen Ausdrucksweisen  neutral(a), Trauer(b) und Freude(c). Der Durchschnittswert der Grundfrequenz beträgt 106 Hz für den neutralen prosodischen Ausdruck (a), 80 Hz für Trauer(b) und 108 Hz für Freude(c).}
	\label{fig:pitch_variation}
\end{figure}




Diese Merkmale beeinflussen die emotionale Sprachproduktion und dienen für den Hörer als akustische Cues für die Emotionswahrnehmung. Bei der emotionalen Sprachproduktion ... 
Studien zeigen eine Tendenz, dass die Grundfrequenz, gefolgt vom Merkmal Sprachtempo
einen starken Einfluss auf die Emotionswahrnehmung haben. \citeauthor{Banse1996} erwähnt eine starke Korrelation zwischen Grundfrequenz und Intensität basierend auf dem physiologischen Prozess der Sprachproduktion: der subglottale Druck führt zur Muskelanspannung, durch die sowohl die Grundfrequenz als auch die Intensität steigen.\\

Aufgrund der Befunde, dass die Grundfrequenz stark mit der Intensität korreliert, wird in dieser Studie der Einfluss untersucht, der eine Änderung der Grundfrequenz oder des Sprachtempos auf die Emotionswahrnehmung in einem interkulturellen Kontext hervorrufen kann.

\subsection{Grundfrequenz}

 BESCHREIBUNG DER MERKMALE! \\

\subsection{Sprachtempo}

Die Wahrnehmung der Variationen des Sprachtempos in der gesprochenen Sprache ist sehr wichtig für die Sprachverarbeitung im Gehirn. Die Änderung/BVariabilität der Aufmerksamkeit der Hörer im Bezug zum Sprachtempo wurde durch perzeptive Experimente belegt (??) \\ 
 

 
Viele Studien zeigen, dass die Hörer auf die Variation der Dauer und Geschwindigkeit der Stimuli reagieren. Die Annahme, dass die Variation der Dauer Auswirkung auf die Emotionswahrnehmung hat, konnte in manchen Studien belegt werden (Scherer 2005?, 2011?). Andere Experimente konnten diese Verbindung nicht eindeutig belegen (\cite{Schirmer2016a})\\

Die Emotionen beeinflussen die Sprachproduktion, da die physiologischen Abläufe, die bei Emotionen entstehen nicht immer unterdrückt oder kontrolliert werden können. Aus der Perspektive der Beschreibungsdimensionen Valenz und Erregung betrachtet, kann die letztere unter Anderem die Frequenz des Herzschlags und die Muskelanspannung beeinflussen. Dies führt zu Änderungen der akustischen Merkmale, die durch den Hörer im Zusammenhang mit den Emotionen verarbeitet werden \parencite{Schirmer2016a}. Damit der Hörer im Stande ist die gehörten Emotionen richtig zu dekodieren, ist zu erwarten, dass die verbalen Ausdrücke eindeutige Hinweise für den Hörer liefern, der dadurch den prosodischen Ausdruck des Sprechers dieser oder jener Emotion zuordnen kann.  In vielen Situationen kommt es dazu, dass diese physiologischen Abläufe ein emotionales Muster im Sprachausdruck hinterlassen. So ist zum Beispiel das Sprachtempo bei Freude eher schneller, während es eher langsamer ist, wenn der Sprecher traurig ist. 
 



1. Datenbereinigung ( Phonotaktik siehe auch unten)\\

Die verwendeten Sprachsignale sind hinsichtlich der Phonotaktik für die deutsche Sprache geprüft, was jedoch für Armenisch nicht gilt.
Nicht alle Phoneme (?) die in Deutsch vorkommen, kommen auch im Armenischen vor.
Deswegen wurden die Aufnahmen hinsichtlich der Phonotaktik des Armenischen (?) geprueft und die Signale, welche in der armenischen Sprache phontaktisch nicht valide (?) sind, aus dem Datensatz herausgenommen. 

Ein Beispiel aus dem Englischen waere, dass im
der schriftlichen Sprache "know" erlaubt ist, also Wörter die mit "kn" anfangen.
In der gesprochenen Sprache wird dieses Wort allerdings wie "no" ausgesprochen, also nicht mit "kn".

Die Schwierigkeit war dabei, dass es ueber die Phonotaktik der gesprochenen Sprache in Armenien nicht ausreichend Literatur gibt. Die geschrieben Sprache wurde zwar ausreichend analysiert, die
gesprochene Sprache dagegen nicht. \\



2. Manipulation der Grundfrequenz\\

Zunächst wurde die Grundfrequenz in den Originalsprachaufnahmen manipuliert. 
Die Änderung der Grundfrequenz wurde nur in eine Richtung durchgeführt - sie wurde erhöht. 
Die Erhöhung der F0-Werte erfolgte in 9 Halbtönen. \\

Exkurs : Grundfrequenz (F0)\\

Die Frequenzen der gesprochenen Sprache einens Menschen sind immer ein Vielfaches der jeweiligen Grundsequenz.
Jeder Mensch verfügt über ein Spektrum von Grundfrequenzen, welches sich im Laufe des Lebens ändert. So hat beispielsweise ein Kind ein Grundsequentspektrum von (?), während Erwachsene im Durchschnitt über ein Spektrum von (?) verfügen.

Die Grundfrequenz entsteht durch die Anspannung der Stimmlippen, welche beim Sprechen durch das Heben und Senken des Kehlkopfes beeinflusst wird..
Angespannte Stimmlippen erzeugen höhere Grundfrequenzen in einem kontinuierlichen Spektrum. So ist die Grundfrequenz bei einem vorderen geschlossen ungerundeten Vokal (z.B. i) ist die
intrinsische Grundfrequenz höher als bei einem vorderen halboffenen ungerundeten Vokal (z.B. e)
Aus der entgegengesetzten Perspektive kann die künstliche Erhöhung der Grundfrequenz zu einer
Änderung in der Vokalqualität führen. Wenn die Grundfrequenz sehr erhöht wird, kann so ein e als i wahrgenommen werden.

Für die Vokalqualität ist zudem die tonotopische Distanz von Bedeutung, also z.b. die
Distanz zwischen f0 und f1.\\

Um zu verhindern, dass bei einer Erhöhung der Grundfrequenz die Vokalqualität leidet, wurden die Originalsignale stufenweise manipuliert. Als Einheit wurde der Halbton benutzt, damit die
Natürlichkeit die Harmonischen bewahrt bleibt, d.h.damit die Stimme natuerlich klingt.
Studien zeigen dass eine Änderung unter 4 Halbtönen für den Menschen nicht hörbar ist. 
Daher sollten die Originalsignale um mehr als 5 Halbtöne verändert werden.\\

Für die Manipulation der Sprachsignale habe ich auf die Pythonbibliothek Parselmouth zugegriffen, die als
Schnittstelle zwischen Python und Praat den Zugriff auf die Praat-Funktionalität ermöglicht (\cite{Jadoul2018a}). Da die Datenvorbereitung für die Originalsprachsignale in einer Pythonumgebung durchgeführt wurde, ist es sinnvoll, die Datenmanipulation in dieser Umgebung vorzunehmen.
Die Erhöhung der Grundfrequenz erfolgte nicht nur bei einzelnen Vokalen, sondern die ganze Grundfrequenzeben timestepp 0,001 sec.wurde extrahiert und alle Werte erhöht. So wurde die komplette Grundfrequenz des Sprachsignals erhöht.\\

Alle 200 Sprachaufnahmen wurden stufenweise um 5,6,7,8 und 9 Halbtöne erhöht, um herauszufinden um wieviel das Sprachsignal erhöht werden kann, damit die höchste Änderung in der Emotionswahrnehmung erreicht wird, ohne dass die Vokalqualität und Natürlichkeit beeinträchtigt wird.\\

Dabei hat sich herausgestellt, dass bei Sprachsignalen deren f0 um 9 Halbtöne erhöht wurden, eine Änderung der Vokalqualität festgestellt werden konnte.
Damit hat sich die Erhöhung um 8 Halbtöne als zielführend erwiesen.\\

Als weiteres Kriterium für Auswahl der manipulierten Sprachsignale wurde die Signalqualtät herangezogen. Diese sollte nicht beeinträchtigt werden.
Dabei erwies sich die Creaky Voice als Ausschlusskriterium.
Bei Originalsignalen, welche mit  mit creaky voice produziert wurden, sind nach der Resynthese der Sprachsignale, Artefakte entstanden. Die Sprachaufnahmen klangen nicht mehr natürlich.
Die Originalsignale mit Creaky Voice wurden aus dem Datensatz herausgenommen.\\

3. Manipulation des Sprachtempos\\

Für die Manipulation des Sprachtempos wurde wiederum die Parcelmouthbibliothek verwendet, um Praat-Befehle auszuführen.
Praat benutzt den Psolar-Algorythmus, d.h das die Sprachsignale werden in kleine Segmente unterteilt. Diese werden dann überlappend übereinanderrgesetzt, um die Dauer des Signals zu verkürzen.
In dieser Arbeit wurde der Psolar-Algorythmus nicht nur auf die Vokale, sondern auf das komplette Sprachsignal angewendet.\\

Dabei war es wichtig, zu beachten, zunächst die Dauerebene des Signals zu extrahieren, sie dann zu verkürzen und anschließend Dauerebene und Sprachsignal zu resyntetisieren, um zu verhindern dass sich die Grundfrequenz durch die Verkürzung der Dauer des Sprachsignals erhöht.\\

Die Manipulation des Sprachtempos wurde ebenfalls stufenweise durchgeführt. Die Dauer wurde um die Faktoren 0,3-0,8 (?)verkürzt.
Die optimale Stufe bedeutet eine hörbare Änderung im Vergleich zum Origanalsignal bei Beibehaltung der Natürlichkeit.
Als bester Faktor stellte sich der Faktor 0,65 heraus.
Eine Verkürzung des Sprachignals um den Faktor 0,65 entspricht einer Erhöhung der Sprachgeschwindigkeit um 35\%




\bigskip

Experiment Design
Experiment als Webexperiment durchgeführt.
Zum designen eines Online-Experimentes bietet sich die Percy-Software (\cite{Draxler2011}) an, da hier die Vorlagen für das Design an die Anforderungen des jeweiligen Experimentes individuell angepasst werden können.


\subsection{Intensität}
\label{intensity}

Intensität als akustischer Merkmal wird durch Amplitude dargestellt.



\section{Interkulturelle Emotionswahrnehmung}


\section{Diskussion}
\label{sec:method_discussion}




\chapter{Emotion und Sprache -  das Experiment}
\label{sec:captain}

Die Emotionswahrnehmung in der gesprochenen Sprache kann durch ein Perzeptionsexperiment evaluiert werden. Die Messkriterien für die Wahrnehmung der Emotionen können entweder kategorial oder numerisch sein. Im letzteren Fall müssen die Hörer die wahrgenommene Emotion auf einer Skala mit einer numerischen Eingabe bewerten. In diesem Fall können auch die Variationen der Perzeption innnerhalb derselben Emotion analysiert werden. Bei der kategorialen Methode müssen die Hörer für den Stimulus einer Kategorie zuordnen. Diese Kategorien können zum Beispiel Emotionen darstellen (Wut, Trauer, Freude usw.).  


Im Rahmen dieser Arbeit wurde ein Web-Experiment in zwei Sprachen vorbereitet - armenisch und deutsch. Die Teilnehmer wurden nach soziodemographischen Daten gefragt, wie Geschlecht, Alter und das Land, in dem sie wohnen. Außerdem gab es Fragen zur Persönlichkeit, die nach dem Fünf-Faktoren-Modell (Big Five) vorbereitet wurden (siehe \ref{appendix:bigfive}). Zu jeder dieser darin enthaltenen Dimensione gab es eine Frage (Offenheit, Verträglichkeit, Extraversion, Gewissenhaftigkeit, Neuroticismus). Die Übersetzung dieser Fragen ins Armenische wurde selbstständig durchgeführt, da keine evaluierte Version des Fragebogens, die für wissenschaftliche Recherche offen zur Verfügung stand, vorhanden war. Die deutsche Version des Fragebogens zur Persönlichkeit ist im Anhang \ref{appendix:bigfive} zu finden. 

Für das Experiment wurde die Percy Software benutzt \autocite{Draxler2011, Draxler2014}. 



Datenschutz: anonym
Meta Daten


Percy
Datenbankstruktur
Anonym


\section{Phonetik und Phonologie der armenischen Sprache}
\label{armenian_phonology}

Silbenstruktur der armenischen Sprache wurde zwar von armenischen Linguisten viel untersucht (\todo{armenische Linguisten zitieren, die über Silbenstruktur geschrieben haben}), allerdings geht es dabei um die Orthography. Dabei wird die gesprochene Sprache nur im Zusammenhang mit den jeweiligen Formen der schriftlichen Sprache diskutiiert. Obwohl das Grapheme-Phonem Verhältnis überwiegend übereinstimmt, nicht alle Buchstaben entsprechen einem Phonem. 

Die Stimuli wurden von den deutschen und armenischen Hörern annotiert. Da die Stimuli für beide Gruppen dieselben Sprachaufnahmen waren, sollten die Pseudowörter, die als Stimuli benutzt wurden, sowohl der deutschen als auch der armenischen Phonotaktik entsprechen. Die Prüfung nach der Legalität der Lautkombinationen für die deutschen Hörer wurde bei der Erstellung des des "MagdeburgerProsodie-Korpus”, aus welchem die Pseudowoerter ausgewählt wurden, durchgeführt (\cite{Wendt2002}). \\

\todo{Was ist uns biher über die Phonotaktik und Phonologie der Armenischen bekannt? H.Hovhannisyan, Vaux, Dum-Tragut.}


Der Mangel an der Literatur zur armenischen Phonotaktik , sowie zur Phonologie und der Phonetik der gesprochenen armenischen Sprache hat dazu geführt, dass die Legalität mancher Lautkombinationen nicht mit Regeln der uns bekannten Literatur geprüft werden konnte. Daher sollten tatsächliche Wörter aus dem Armenischen gefunden werden, die die Lautfolgen der Pseudowörter enthalten. 


Alle Phoneme, die in einer Sprache vorkommen werden im Phoneminventar der jeweiligen Sprache aufgeführt. Jedoch sind nicht alle Kombinationen dieser Laute legal. Da Hörer illegale Lautkombinationen als unnatürlich empfinden,  sollten diese in Stimulis nicht vorkommen.\\


Armenisches Phoneminventar unterscheidet sich 



\section{Motivation}
\label{sec:experimemnt_motivation}

\section{Methodik}
\label{sec:method}

Für das in dieser Arbeit beschriebene Perzeptionsexperiment wurden Stimuli benötigt, welche durch den Hörer evaluiert werden können. Dazu wurden die Merkmale der Originalsprachaufnahmen künstlich verändert. Es stellt sich die Frage, ob eine Änderung der Stimuli Einfluss auf die Emotionswahrnehmung hat.\\


Um dies zu untersuchen sind mehrere unterschiedliche Methoden denkbar, die sich sowohl in ihrer Aussagekraft als auch in ihrer Durchführbarkeit unterscheiden.\\

Bei einer Methode werden die manipulierten Stimuli anhand vorgegebener Kategorien, wie zum Beispiel der Intensität einer Emotion, auf einer Skala bewertet. Da diese Arbeit aber das Ziel hat, die Wahrnehmung der Emotionen zu untersuchen, würde dies durch diesen Ansatz gerade durch die Vorgabe der Kategorien verhindert werden.\\

Eine andere Methode würde sowohl die Auswahl der Kategorie als auch die Bewertung auf einer Skala dem Probanden überlassen. Durch diesen Anstieg an Freiheitsgraden stiege aber auch die Dimensionalität der Daten, wodurch auch die Anzahl der Stichproben drastisch gesteigert werden müsste, um eine statistische Signifikanz zu erreichen. Dies würde den Rahmen dieser Arbeit übersteigen.\\

Deshalb wurde eine Methode ausgewählt, die Stimuli nicht willkürlich sondern gezielt manipuliert und diese dann vom Probanden einer vorgegebenen Liste von Emotionskategorien zugeordnet werden soll. Dadurch wird sichergestellt, dass der Hörer die Entscheidung trifft, welcher Emotion er das Sprachsignal zuordnet. Auf der anderen Seite bleibt das Experiment relativ übersichtlich und mit der realistisch erwartbaren Stichprobenpopulation statistisch auswertbar.

\subsection{Auswahl der Merkmale}

\subsection{Auswahl des Sprachkorpus}
\label{sec:db}

Für die Untersuchung der interkulturellen Emotionswahrnehmung musste EMotionswahrnehmung im Mittelpunkt stehen. Daher sollten möglichst wenige Störfaktoren vorhanden sein. Ein sehr wichtiger Störfaktor für die Untersuchung der Emotionswahrnehmung ist die Verarbeitung dr semantischen Aspekte der Sprachperzeption im Gehirn, da die Semantik ebenfalls Emotionen auslösen kann. Aus diesem Grund sollte die Prosodie der Stimuli der Hauptauslöser für die Emotionen sein (eigentlich sollten keine Emotionen ausgelöst werden, sondern wahrgenommen). 

Der Magdeburger Prosodie-Korpus besteht aus zwei Sprachdatenbanken. Die eine enthält deutsche Substantive, die zweite besteht aus Pseudowörtern. Beide Teilen wurden von einem Mann und von einer Frau ausgesprochen. Die sprecher waren professionelle Schauspeler und haben alle Substantive und Pseudo-Wörter in sechs prosodischen Ausdrücke ausgesprochen. Durch Perzeptionsexperimente wurde die die Akzeptanz der Emotionen in Sprachaufnahmen untersucht \autocite{Wendt2003}.

Die Validierung der ausgesprochenen Emotionen in den Sprachaufnahmen wurde durch Perzeptionsexperimente mit trainierten sowie mit untrainierten Hörern durchgeführt. Die Ergebnisse der Evaluation der prosodischen Ausdrücke weisen eine Akzeptanz der Emotionen über 70 Prozent auf. Eine Ausnahme ist der prosodischer Audruck Trauer in Pseudowörtern, die von dem männlichen Schauspieler ausgesprochen wurden (\ref{fig:wasep_evaluation}).

\begin{figure}
	\centering
	\includegraphics[width=1\linewidth]{pics/wasep_evaluation.png}
	\caption{\textbf{Die Akzeptanz der ausgesprochenen Emotionen bei allen Hörern (aus \cite{Wendt2002}).}}
	\label{fig:wasep_evaluation}
\end{figure}

beinhaltet Sprachaufnahmen in folgenden Basisemotionen: Freude, Wut, Angst, Trauer, Ekel sowie neutral. 

Für das Experiment wurde zuerst Recherchiert und einen Sprachkorpus ausgesucht, der "Der Magdeburger Prosodie Korpus " heisst: der Magdeburger Korpus für Prosodie. Dabei wurde von Dr. Wendt untersucht bei durch fMRI Bildern wurde aufgezeichnet die emotionale Sprache und dann hat sie mehrere Experimente durchgeführt und 200 Sprachdatenbanken daraus gemacht. Eine Datenbank enthält deutsche Wörter, die zweite Datenbank enthält Pseudowörter. Diese Pseudowörter wurden aus deutschen Wörtern erstellt, indem die Buchstaben umsortiert wurden und daraus hat man Pseudowörter erhalten, die der deutschen Phonotaktik entsprechen. Diese Pseudowörter wurden von 2 deutschen Schauspielern , ein Mann und eine Frau in unterschiedlichen Emotionen ausgesprochen. Diese Basisemotionen sind, Trauer, Wut, Freude, Angst, Ekel und neutral, also es sind 5+1 Emotionen. Jedes Wort wurde sowohl von dem SPrecher als auch von der Sprecherin in allen vorgegebenen Basisemotionen sowie neutral ausgesprochen.


\section{Pseudowörter}

Es handelt sich dabei um 200 Pseudowörter. Es wurde entschieden Stimuli aus diesen Pseudowörter auszuwählen, da diese keine Bedeutung weder für die Deutschen noch für Armenier haben. Also können die Hörer durch die semantische Bedeutung nicht abgekelenkt werden, und dadurch kann möglichst mehr Aufmerksamkeit in der Verarbeitung im Gehirn auf Emotion gelegt werden. Aus diesen Pseudowörtern aurden dann Stimulis folgenderweise ausgesucht. Alle 200 Ausfnahmen wurden zunächst nach der armenischen Phonotaktit geprüft, und viele Wörter ausgefallen sind. Zum Beispiel gab es manche Wörter, die einem armenischen Wort ähnlich waren.  Bei vielen Wörtern würde das nicht in Frage komme, da die Wörter, die mit Vokal anfangen, werden von den Deutschen mit Glottal Stop ausgesprochen. Auf Armenisch gibt es kein Glottal Stop vor einem Vokal am Wortanfang. Daher wurden alle Wörter herausgefiltert, die mit einem Vokal anfangen. Dann kommen die Wörter, die die entsprechenden Phoneme auf Armenisch nicht haben. Alle Pseudowörter, die Laute enthalten, die nicht im armenischen Phoneminventar vorkommen oder nicht als allophone eines Phonems wahrgenommen werden konnten, wurden von der Auswahl ausgeschlossen. Und dann wurden alle Wörter die Klaster haben, die im Armenischen nicht legal sind, ebenfallls herausgefiltert. Am Ende waren es 15 Wörter, die der armenischen Phonotaktik entsprechen. Aus diesen 15 Wörtern wurden demnächst Stimuli vorbereitet. Also die Manipulationen für F0 und Tempo wurden durchgeführt, und nach einem ernueten Prüfen von Stimuli, wurden die gelöscht, die durch schlechte Qualität Artefakte hatten. Am Ende sind 4 Stimuli gebblieben, die für Experiment in Frage kommen. Die Artefakte waren Verzehrungen der Stimme bei einer Emotion eines Sprechers, und das würde genügen, dass dieser Stimulus nicht miteinbezogen werden konnte.

\section{Stimuli}
\label{sec:captain_overview}

Die Emotionswahrnehmung in Sprache kann durch die semantische Bedeutung beeinflusst werden. Da die Pseudowörter keine Bedeutung auf Deutsch haben, waren diese als Stimuli geeignet. Es befanden sich 200 Pseudowörter in dem Sprachkorpus, die zunächst detailliert untersucht wurden, um daraus die Stimuli auszuwählen. Für diese Auswahl mussten die Pseudowörter folgende Bedingungen erfüllen:

\begin{itemize}
	\item das Wort hat keine Bedeutung auf Armenisch
	\item es enthält nur Phoneme, die im Armenischen als Phoneme oder als Allophone vorkommen
	\item es entspricht der armenischen Phonotaktik
	\item nach der Manipulation des Sprachsignals entstehen keine Artefakte 
\end{itemize} 




\subsection{Wortauswahl}
\label{sec:pseudowords}

Für die Auswahl der Wörter wurden folgende Ausschlusskriterien festgelegt, wobei die Kriterien eins bis vier bei der Vorauswahl eine Rolle spielten, wohingegen Kriterium fünf erst bei der Nachbearbeitung zum
Einsatz kam.


\begin{enumerate}
	\item Das Pseudowort enthält	mindestens ein Phonem, das im armenischen Phoneminventar nicht vorkommt.
	\item Das Pseudowort entspricht nicht der armenischen Phonotaktik, das heisst es beinhaltet mindestens eine
	illegale Lautkombination für die armenischen Hörer.
	\item Das Pseudowort beginnt mit einem Vokal.
	Da die Sprecher der Pseudowörter Deutsche sind, werden Wörter die mit einem Vokal beginnen
	mit einem glottal stop ausgesprochen, was bedeutet, dass der Onset der ersten Silbe ein glottal Stop ist. 
	Im Armenischen hingegen wird nicht unbedingt ein glottal stop benutzt, daher könnte dies als unnatürlich empfunden werden und deshalb wurden diese Wörter aussortiert.
	\item Die Aussprache des Pseudowortes zeigt auffällige Unterschiede, wenn es von verschiedenen Sprechern ausgesprochen wird. 
	\item Mindestens ein manipuliertes Sprachsignal des Pseudowortes entspricht nicht mehr den geforderten
	Qualitätsansprüchen an die Natürlichkeit, zum Beispiel weil es bei der Resynthese zu einer Artefaktbildung kam.	
\end{enumerate}


(Bild Beispiel mit Artefakten)
Caption: Bei manchen Sprachaufnahmen waren Artefakte nach der Resynthese zu hören. Diese Aufnahmen wurden aus der Stimuli-Liste herausgenommen.
(Bild: armenisches Phoneminventar vielleicht im Anhang)\\

Da die verfügbare Information über die armenische Phonotaktik in der Literatur für diese Überprüfung nicht ausreichend war, basiert sie auf Beispielwörtern aus dem Armenischen, in denen die Lautkombinationen der Pseudowörter vorkommen.

Manche Autoren, die die armenische Phonotaktik dargestellt haben, stellen die orthographischen Regeln dar (\todo{Ajarian missing, andere Autoren für Phonetik missing, und die dum-targuto.ä. }).



{
	\renewcommand{\arraystretch}{1.2}
	\begin{table}[!t]
		\centering
		
		\caption{\textbf{Die folgenden Pseudowörter wurden als Stimuli ausgewählt.} Nach der Prüfung der Legalität für armenische Hörer, sowie der Qualitätsprüfung sind nur 4 Pseudowörter zur Auswahl geblieben.} \label{table:stimuli}
		
		\begin{tabular}{@{}lllll}
			
			\toprule
			
			\textbf{PSEUDOWORT} & \textbf{KAN} & \textbf{IPA} & \textbf{ARM IPA} & \textbf{ARM}\\
			
			\midrule
			
			tápik & ta:pIk & \textipa{[tAp\super{h}ik\super{h} ]} & \textipa{[tatik] [mArt\super{h}ik\super{h}]} & {\artm tatik, mardik } \\
			misuk & mizUk &	\textipa{[mizuk\super{h}]} & \textipa{[mitum] [bAzuk] [mORuk\super{h}]} & {\artm mijaket, bazuk, moruq } \\			
			sawet &	zav@t & \textipa{[zAvEt\super{h}]} & \textipa{[zAtik] [gEt\super{h}]} & {\artm Zaven, havet, geth } \\
			nemmok & nemOk & \textipa{[nEmOk\super{h}]} & \textipa{[nEmRut\super{h}] [kAnEm] [kuzEs]} & {\artm Nemruth, kanem, kuzes, kamoq} \\		
			
			
			
			\bottomrule
			
		\end{tabular}
		
	\end{table}
}



\subsection{Werkzeuge: eigene und bestehende}

\subsection{Manipulation der Sprachsignale}

\subsubsection{Manipulation der Grundfrequenz}
\label{sec:f0_manipulation}
F0 steht im Verhältnis zur Vokalhöhe (Peterson and Barney, 1952; Ohala, 1937; Traunmüller, 1995)

Parselmouth ist eine Schnittstelle zwischen Python und Praat. Durch Parselmouth können Praat befehle aus einem Python-Skript durchgeführt werden. Die Manipulation von Stimuli wurden mittels Parselmouth durchgeführt. 		
F0 

wenn F0-F1 Abstand viel zu gering ist, ändert sich die Vokalqualität (front mid-open -> front mid-closed)
daher soll die Pitch-Änderung nicht größer sein als 6 Halbtöne, aber auch nicht kleiner als 5 Halbtöne.

- F0 Manipulator : 
Parameter "Shift Frequency":
Time-Step: 0.001
Unit: semitones



\begin{figure}
	\centering
	\begin{subfigure}{1\textwidth}
		\centering
		\includegraphics[width=1\linewidth]{pics/nemmok.png}
		\caption{original}
		\label{fig:nemmok_original}
	\end{subfigure}
	\par\bigskip
	\begin{subfigure}{1\textwidth}
		\centering
		\includegraphics[width=1\linewidth]{pics/nemmok_pitch.png}
		\caption{f0 manipuliert}
		\label{fig:nemmok_pitch}
	\end{subfigure}
	\caption{\textbf{Stimulusvergleich: f0 unverändert(a) und manipuliert(b) .} Oszillogramm(oben) und Spektrogramm (unten) für das Pseudowort \emph{nemmok} (IPA: \textipa{[nEmOk\super{h}]}). Prosodischer Ausdruck neutral augesprochen mit Frauenstimme. Die blaue Linie auf dem Spekrogramm zeigt den Verlauf der Grundfrequenz. Die Grundfrequenzwerte im gesamten Originalsprachsignal (a) wurden um 8 Halbtöne erhöht (b). }
	\label{fig:f0_manipulation}
\end{figure}


Die Grundfrequenz wurde zunächst stufenweise (2, 3, 4, 5, 6, 7, 8 und 9 Halbtöne) erhöht. Anschließend wurden die manipulierten Sprachsignale mit den originalen Sprachaufnahmen verglichen. Dadurch konnte es festgestellt werden, ob die Änderung der Grundfrequenz durch das menschliche Gehör wahrgenommen werden kann. \todo{Es wird vermutet, dass die Änderung unter 5 Halbtönen nicht als Änderung wahrgenommen kann ???? REFERENZ FINDEN DRINGEND! Burkhardt? oder Zitat bei Burkhardt? sonst 'pitch shift emotion perception' suchen}\\

Eine Pitch-Änderung um weniger als 4 Halbtöne macht keinen großen perzeptiven Unterschied. 
Entgültige Entscheidung: eine Pitch-Änderung um 8 Halbtöne \\

Achtung: bei creaky voice entstehen unschöne Artefakte. Also wurden alle Sprachaufnahmen, die creaky voice enthalten, aus der weiteren Auswahl ausgeschlossen.

\subsubsection{Manipulation des Sprachtempos}
\label{sec:tempo}

Die Dauer wurde zunächst stufenweise um Faktor 0.5, 0.55, 0.6, 0.65, 0.7, 0.75, 0.8 und 0.9 skaliert. Anschließend wurde Faktor 0.65 ausgewählt, da die Sprachsignale einen Unterschied zum Originalaufnahmen aufweisen.

Parameter "Scale times by":
Factor: 0.65

Verkürzung der Dauer um Faktor 0.65 heisst 35 Prozent schneller.


\begin{figure}
	\centering
	\begin{subfigure}{1\textwidth}
		\centering
		\includegraphics[width=1\linewidth]{pics/nemmok.png}
		\caption{original}
		\label{fig:nemmok_original}
	\end{subfigure}
	\par\bigskip  
	\begin{subfigure}{1\textwidth}
		\centering
		\includegraphics[width=1\linewidth]{pics/nemmok_tempo.png}
		\caption{Dauer manipuliert}
		\label{fig:nemmok_tempo}
	\end{subfigure}
	\caption{\textbf{Stimulusvergleich: Dauer unverändert(a) und manipuliert(b) .} Oszillogramm(oben) und Spektrogramm (unten) für das Pseudowort \emph{nemmok} (IPA: \textipa{[nEmOk\super{h}]}). Prosodischer Ausdruck neutral augesprochen mit Frauenstimme. Die Dauer des Originalsprachsignals beträgt ~0.96 s. (a), nach der Manipulation der Dauer um Faktor 0.65 beträgt die Dauer des Sprachsignals (b) ~0.62 s. Dies entspricht einer Erhöhung der Sprechgeschwindigkeit um 35\%. Die blaue Linie auf dem Spekrogramm zeigt den Verlauf der Grundfrequenz. Bei der Manipulation der Dauer ist die Grundfrequenz gleich geblieben.}
	\label{fig:tempo_manipulation}
\end{figure}


\subsection{Experimentdesign}
\label{sec:experiment_design}

Das Experiment wurde als eine explanative Feldstudie durchgeführt. 
\newline


Die verfügbaren Stimuli 

{
	\renewcommand{\arraystretch}{1.2}
	\begin{table}[!t]
		\centering
		
		\caption{\textbf{Nach der Manipulation der Sprachsignale hat sich für jedes Experiment folgendes Bild ergeben:}} \label{table:design}
		
		\begin{tabular}{@{}lll}
			
			\toprule
			
			\textbf{Original} & \textbf{F0} & \textbf{Tempo}\\
			
			\midrule
			
			4 Wörter * 2 Sprecher * 6 prosodische Ausdrücke & 
			4 Wörter * 2 Sprecher * 6 prosodische Ausdrücke & 4 Wörter * 2 Sprecher * 6 prosodische Ausdrücke\\
			
			\midrule
			
			48 Sprachsignale & 48 Sprachsignale & 48 Sprachsignale \\	
			
			\bottomrule
			
			\textbf{insgesamt 144 Stimuli} \\				
			
		\end{tabular}
		
	\end{table}
}

Mit Hilfe der Schneeball-Methode wurde der Link zu diesem Web-Experiment verbreitet. 
Für die qualitätsprüfung der Annotationen wurde eine Kontrollfrage eingebaut.
Die Kontrollfrage sollte dazu dienen, die Annotationen der Teilnehmer, die nicht Aufmerksam waren und zufällige Emotionen für die Stimuli gewählt haben, aus der weiteren Datenanalyse auszuschließen. Ein Audioaufnahme mit Applaus wurde als Stimulus für die Kontrollfrage benutzt. Dazu sollten die Teilnehmer die Option 'Klatschen' anstatt einer Emotion wählen. Es wurden nur die Annotationen der Höher analysiert, die die Kontrollfrage richtig beantwortet haben. Außerdem wurden die Daten für jeden Teilnehmer hinsichtlich folgender Kriterien geprüft: 

\begin{itemize}
	\item Muttersprache der Teilnehmer des deutschen Experimentes ist Deutsch 
	\item Muttersprache der Teilnehmer des armenischen Experimentes ist Armenisch
	\item Deutsch kommt in der Liste der angegebenen Fremdsprachen nicht vor
	\item Armenisch kommt in der Liste der angegebenen Fremdsprachen nicht vor
	\item es sind 151 Einträge für den Teilnehmer vorhanden, dabei versteht es sich folgendes:
	
	\begin{itemize}	 
		\item Annotationen für alle Stimuli (insgesamt 144 Sprachaufnahmen und eine Kontrollfrage) sind vorhanden
		\item alle Persönlichkeitsfragen (insgesamt fünf Fragen) wurden beantwortet
		\item die Frage nach der Natürlichkeit der Stimme von den Sprechern wurde beantwortet
	\end{itemize}
\end{itemize}



Das Experiment sollte möglichst erreichbar für die Teilnehemer der beiden zu vergleichenden Länder sein. Außerdem sollten in eher kuzer Zeit möglichst viele Personen an dem Experiment teilnehmen. Für diese Ziele hat sich ein Web-Experiment sehr gut geignet, das die Erreichbarkeit für ein solches Experiment höher ist, als für ein Laborexperiment, bei dem die Teilnehmer persönlich und wahrscheinlich einzeln eintreten müssten. Daher wurde es entschieden, dass sich ein Web-Experiment am besten für die aktuelle Studie eignet. 

Die verfügbaren Platformen, die benutzt werden konnten, um ein online Experiment durchzuführen, wurden  zunächst verglichen. Percy Software \parencite{Draxler2011} 


\subsection{Percy - die Praesentation der Stimuli und Fragen fuer die Hoerer}
\label{sec:stimulusobject}

Die Teilnehmer hatten die Aufgabe alle Stimuli zu kategorisieren, indem sie für jedes Sprachsignal eine Emotion wählen. 
\begin{figure}
	\centering
	\includegraphics[width=1\linewidth]{pics/experiment/experiment_active_buttons.png}
	\caption{\textbf{Diskriminationsaufgabe.} Die Frage zur Kategorisierung des Stimulus zu einer Emotion. }
	\label{fig:experiment_question}
\end{figure}


 Der Stimulus konnte mit dem Kopfhörersymbol abgespielt werden. Für alle Stimuli wurde den Teilnehmern die selbe Frage gestellt. Für die Aktivierung der Auswahloptionen musste der Stimulus mindestens einmal abgespielt werden. Die Teilnehmer durften jede Sprachaufnahme maximal dreimal hören. Sobald der Teilnehmer eine der unter der Frage dargestellten Optionen ausgewählt hat, wurde der nächste Stimulus präsntiert. Es konnte nur eine der angegebenen Optionen ausgewählt werden.
 
\subsubsection{Fragen zur Persönlichkeit}

\subsubsection{Kontrollfrage}

\subsection{Probandenauswahl a posteriori}


 
\chapter{Datenanalyse}

\section{Demographische Auswertung der Teilnehmerdaten}

wie viele Männer/Frauen haben bei jedem experiment teilgenommen? 
Durschnittsalter und Standardabweichung


\section{Vergleich der eigenen Daten mit der Evaluation des Magdeburger Sprachkorpus}

\section{Vergleich der Emotionswahrnehmung von armenischen und deutschen Hörern}

\begin{figure}[t!] % "[t!]" placement specifier just for this example
	\begin{subfigure}{0.48\textwidth}
		\includegraphics[width=\linewidth]{plots/heatmap/conf_original_de.png}
		\caption{First subfigure} \label{fig:conf_original_de}
	\end{subfigure}\hspace*{\fill}
	\begin{subfigure}{0.48\textwidth}
		\includegraphics[width=\linewidth]{plots/heatmap/conf_original_am.png}
		\caption{Second subfigure} \label{fig:conf_original_de}
	\end{subfigure}
	
	\medskip
	\begin{subfigure}{0.48\textwidth}
		\includegraphics[width=\linewidth]{plots/heatmap/conf_pitch_de.png}
		\caption{Third subfigure} \label{fig:conf_pitch_de}
	\end{subfigure}\hspace*{\fill}
	\begin{subfigure}{0.48\textwidth}
		\includegraphics[width=\linewidth]{plots/heatmap/conf_pitch_am.png}
		\caption{Fourth subfigure} \label{fig:conf_pitch_am}
	\end{subfigure}
	
	\medskip
	\begin{subfigure}{0.48\textwidth}
		\includegraphics[width=\linewidth]{plots/heatmap/conf_tempo_de.png}
		\caption{Fifth subfigure} \label{fig:conf_tempo_de}
	\end{subfigure}\hspace*{\fill}
	\begin{subfigure}{0.48\textwidth}
		\includegraphics[width=\linewidth]{plots/heatmap/conf_tempo_am.png}
		\caption{Sixth subfigure} \label{fig:conf_tempo_am}
	\end{subfigure}
	
	\caption{My complicated figure} \label{fig:confusion_heatmap}
\end{figure}



\subsection{intrakulturelle Unterschiede, Persoenlichkeitsunterschied}

\subsection{Freude}


\begin{figure}[t!] % "[t!]" placement specifier just for this example
	\begin{subfigure}{0.48\textwidth}
		\includegraphics[width=\linewidth]{plots/de_Freude.png}
		\caption{First subfigure} \label{fig:de_F}
	\end{subfigure}\hspace*{\fill}
	\begin{subfigure}{0.48\textwidth}
		\includegraphics[width=\linewidth]{plots/am_Freude.png}
		\caption{Second subfigure} \label{fig:am_F}
	\end{subfigure}
	
	\medskip
	\begin{subfigure}{0.48\textwidth}
		\includegraphics[width=\linewidth]{plots/de_Freude_F0.png}
		\caption{Third subfigure} \label{fig:de_F_F0}
	\end{subfigure}\hspace*{\fill}
	\begin{subfigure}{0.48\textwidth}
		\includegraphics[width=\linewidth]{plots/am_Freude_F0.png}
		\caption{Fourth subfigure} \label{fig:am_F_F0}
	\end{subfigure}
	
	\medskip
	\begin{subfigure}{0.48\textwidth}
		\includegraphics[width=\linewidth]{plots/de_Freude_tempo.png}
		\caption{Fifth subfigure} \label{fig:de_F_tempo}
	\end{subfigure}\hspace*{\fill}
	\begin{subfigure}{0.48\textwidth}
		\includegraphics[width=\linewidth]{plots/am_Freude_tempo.png}
		\caption{Sixth subfigure} \label{fig:am__F_tempo}
	\end{subfigure}
	
	\caption{My complicated figure} \label{fig:countplots_F}
\end{figure}

\subsection{Trauer}

\begin{figure}[t!] % "[t!]" placement specifier just for this example
	\begin{subfigure}{0.48\textwidth}
		\includegraphics[width=\linewidth]{plots/de_Trauer.png}
		\caption{First subfigure} \label{fig:de_T}
	\end{subfigure}\hspace*{\fill}
	\begin{subfigure}{0.48\textwidth}
		\includegraphics[width=\linewidth]{plots/am_Trauer.png}
		\caption{Second subfigure} \label{fig:am_T}
	\end{subfigure}
	
	\medskip
	\begin{subfigure}{0.48\textwidth}
		\includegraphics[width=\linewidth]{plots/de_Trauer_F0.png}
		\caption{Third subfigure} \label{fig:de_T_F0}
	\end{subfigure}\hspace*{\fill}
	\begin{subfigure}{0.48\textwidth}
		\includegraphics[width=\linewidth]{plots/am_Trauer_F0.png}
		\caption{Fourth subfigure} \label{fig:am_T_F0}
	\end{subfigure}
	
	\medskip
	\begin{subfigure}{0.48\textwidth}
		\includegraphics[width=\linewidth]{plots/de_Trauer_tempo.png}
		\caption{Fifth subfigure} \label{fig:de_T_tempo}
	\end{subfigure}\hspace*{\fill}
	\begin{subfigure}{0.48\textwidth}
		\includegraphics[width=\linewidth]{plots/am_Trauer_tempo.png}
		\caption{Sixth subfigure} \label{fig:am__T_tempo}
	\end{subfigure}
	
	\caption{My complicated figure} \label{fig:countplots_T}
\end{figure}


\subsection{Neutral}

\begin{figure}[t!] % "[t!]" placement specifier just for this example
	\begin{subfigure}{0.48\textwidth}
		\includegraphics[width=\linewidth]{plots/de_neutral.png}
		\caption{First subfigure} \label{fig:de_s}
	\end{subfigure}\hspace*{\fill}
	\begin{subfigure}{0.48\textwidth}
		\includegraphics[width=\linewidth]{plots/am_neutral.png}
		\caption{Second subfigure} \label{fig:am_s}
	\end{subfigure}
	
	\medskip
	\begin{subfigure}{0.48\textwidth}
		\includegraphics[width=\linewidth]{plots/de_neutral_F0.png}
		\caption{Third subfigure} \label{fig:de_s_F0}
	\end{subfigure}\hspace*{\fill}
	\begin{subfigure}{0.48\textwidth}
		\includegraphics[width=\linewidth]{plots/am_neutral_F0.png}
		\caption{Fourth subfigure} \label{fig:am_s_F0}
	\end{subfigure}
	
	\medskip
	\begin{subfigure}{0.48\textwidth}
		\includegraphics[width=\linewidth]{plots/de_neutral_tempo.png}
		\caption{Fifth subfigure} \label{fig:de_s_tempo}
	\end{subfigure}\hspace*{\fill}
	\begin{subfigure}{0.48\textwidth}
		\includegraphics[width=\linewidth]{plots/am_neutral_tempo.png}
		\caption{Sixth subfigure} \label{fig:am__s_tempo}
	\end{subfigure}
	
	\caption{My complicated figure} \label{fig:countplots}
\end{figure}


\subsection{Ekel}

\begin{figure}[t!] % "[t!]" placement specifier just for this example
	\begin{subfigure}{0.48\textwidth}
		\includegraphics[width=\linewidth]{plots/de_Ekel.png}
		\caption{First subfigure} \label{fig:de_E}
	\end{subfigure}\hspace*{\fill}
	\begin{subfigure}{0.48\textwidth}
		\includegraphics[width=\linewidth]{plots/am_Ekel.png}
		\caption{Second subfigure} \label{fig:am_E}
	\end{subfigure}
	
	\medskip
	\begin{subfigure}{0.48\textwidth}
		\includegraphics[width=\linewidth]{plots/de_Ekel_F0.png}
		\caption{Third subfigure} \label{fig:de_E_F0}
	\end{subfigure}\hspace*{\fill}
	\begin{subfigure}{0.48\textwidth}
		\includegraphics[width=\linewidth]{plots/am_Ekel_F0.png}
		\caption{Fourth subfigure} \label{fig:am_E_F0}
	\end{subfigure}
	
	\medskip
	\begin{subfigure}{0.48\textwidth}
		\includegraphics[width=\linewidth]{plots/de_Ekel_tempo.png}
		\caption{Fifth subfigure} \label{fig:de_E_tempo}
	\end{subfigure}\hspace*{\fill}
	\begin{subfigure}{0.48\textwidth}
		\includegraphics[width=\linewidth]{plots/am_Ekel_tempo.png}
		\caption{Sixth subfigure} \label{fig:am__E_tempo}
	\end{subfigure}
	
	\caption{My complicated figure} \label{fig:countplots_E}
\end{figure}


\subsection{Angst}
\begin{figure}[t!] % "[t!]" placement specifier just for this example
	\begin{subfigure}{0.48\textwidth}
		\includegraphics[width=\linewidth]{plots/de_Angst.png}
		\caption{First subfigure} \label{fig:de_A}
	\end{subfigure}\hspace*{\fill}
	\begin{subfigure}{0.48\textwidth}
		\includegraphics[width=\linewidth]{plots/am_Angst.png}
		\caption{Second subfigure} \label{fig:am_A}
	\end{subfigure}
	
	\medskip
	\begin{subfigure}{0.48\textwidth}
		\includegraphics[width=\linewidth]{plots/de_Angst_F0.png}
		\caption{Third subfigure} \label{fig:de_A_F0}
	\end{subfigure}\hspace*{\fill}
	\begin{subfigure}{0.48\textwidth}
		\includegraphics[width=\linewidth]{plots/am_Angst_F0.png}
		\caption{Fourth subfigure} \label{fig:am_A_F0}
	\end{subfigure}
	
	\medskip
	\begin{subfigure}{0.48\textwidth}
		\includegraphics[width=\linewidth]{plots/de_Angst_tempo.png}
		\caption{Fifth subfigure} \label{fig:de_A_tempo}
	\end{subfigure}\hspace*{\fill}
	\begin{subfigure}{0.48\textwidth}
		\includegraphics[width=\linewidth]{plots/am_Angst_tempo.png}
		\caption{Sixth subfigure} \label{fig:am__A_tempo}
	\end{subfigure}
	
	\caption{My complicated figure} \label{fig:countplots_A}
\end{figure}


\subsection{Wut}

\begin{figure}[t!] % "[t!]" placement specifier just for this example
	\begin{subfigure}{0.48\textwidth}
		\includegraphics[width=\linewidth]{plots/de_Wut.png}
		\caption{First subfigure} \label{fig:de_W}
	\end{subfigure}\hspace*{\fill}
	\begin{subfigure}{0.48\textwidth}
		\includegraphics[width=\linewidth]{plots/am_Wut.png}
		\caption{Second subfigure} \label{fig:am_W}
	\end{subfigure}
	
	\medskip
	\begin{subfigure}{0.48\textwidth}
		\includegraphics[width=\linewidth]{plots/de_Wut_F0.png}
		\caption{Third subfigure} \label{fig:de_W_F0}
	\end{subfigure}\hspace*{\fill}
	\begin{subfigure}{0.48\textwidth}
		\includegraphics[width=\linewidth]{plots/am_Wut_F0.png}
		\caption{Fourth subfigure} \label{fig:am_W_F0}
	\end{subfigure}
	
	\medskip
	\begin{subfigure}{0.48\textwidth}
		\includegraphics[width=\linewidth]{plots/de_Wut_tempo.png}
		\caption{Fifth subfigure} \label{fig:de_W_tempo}
	\end{subfigure}\hspace*{\fill}
	\begin{subfigure}{0.48\textwidth}
		\includegraphics[width=\linewidth]{plots/am_Wut_tempo.png}
		\caption{Sixth subfigure} \label{fig:am__W_tempo}
	\end{subfigure}
	
	\caption{My complicated figure} \label{fig:countplots_W}
\end{figure}


\chapter{Fazit}



%\printbibliography[title=References]
\printbibliography



\appendix

% No section numbering at any depth ...
\setcounter{secnumdepth}{-1}

% ... such that this chapter is not numbered (like, for instance, the References chapter)
\chapter{Appendix}

% Restore section numbering (2 is the default and means that sections and subsections are numbered, but subsubsections are not)
\setcounter{secnumdepth}{2}

% Configure how sections in the appendix are numbered
\renewcommand\thesection{\Alph{section}}

\section{Fragen zur Persönlichkeit}
\label{appendix:bigfive}

Diese sind fünf Fragen zur Person nach dem Fünf-Faktoren-Model.   

\begin{itemize}
	\item Wie offen und aufgeschlossen sind Sie für neue Erfahrungen?
	\begin{enumerate}
		\item konservativ/vorsichtig
		\item etwas konservativ/vorsichtig
		\item neutral
		\item etwas erfinderisch/neugierig
		\item erfinderisch/neugierig
	\end{enumerate}
	\item Wie gewissenhaft sind Sie?
	\begin{enumerate}
		\item unbekümmert/nachlässig
		\item etwas unbekümmert/nachlässig
		\item neutral
		\item etwas effektiv/organisiert
		\item effektiv/organisiert
	\end{enumerate}
	\item Wie extrovertiert schätzen Sie sich ein?
	\begin{enumerate}
		\item zurückhaltend/reserviert
		\item etwas zurückhaltend/reserviert
		\item neutral
		\item etwas gesellig/aktiv
		\item gesellig/aktiv
	\end{enumerate}
	\item Wie verträglich schätzen Sie sich ein?	
	\begin{enumerate}
		\item wettbewerbsorientiert/antagonistisch
		\item etwas wettbewerbsorientiert/antagonistisch
		\item neutral
		\item etwas kooperativ/mitfühlend
		\item kooperativ/mitfühlend
	\end{enumerate}
	\item Wie impulsiv sind Sie?	
	\begin{enumerate}
		\item selbstsicher/ruhig
		\item etwas selbstsicher/ruhig
		\item neutral
		\item etwas emotional/mitfühlend
		\item emotional/mitfühlend
	\end{enumerate}
\end{itemize}


\noindent Die oben genannten fünf Fragen zur Persönlichkeit wurde für die armenische Version des Experimentes ins Armenische übersetzt.  

\begin{itemize}
	{\artm 
	\item Ինչպե՞ս եք վերաբերվում փոփոխություններին։
	\begin{enumerate}
		\item զգուշավոր/փոփոխամերժ
		\item որոշ չափով զգուշավոր/փոփոխամերժ
		\item չեզոք
		\item որոշ չափով հեշտ հարմարվող/հետաքրքրասեր
		\item հեշտ հարմարվող/հետաքրքրասեր
	\end{enumerate}
	\item Պարտականությունները կատարելիս հիմնականում միտված եք լինելու․․․
	\begin{enumerate}
		\item անհոգ/ոչ հետևողական		
		\item որոշ չափով անհոգ/ոչ հետևողական
		\item չեզոք
		\item որոշ չափով պարտաճանաչ/կազմակերպված
		\item պարտաճանաչ/կազմակերպված
	\end{enumerate}
	\item Համարում եք Ձեզ ավելի շփվո՞ղ, թե՞ ինքնամփոփ մարդ։
	\begin{enumerate}
		\item ինքնամփոփ/զուսպ
		\item որոշ չափով ինքնամփոփ/զուսպ
		\item չեզոք
		\item որոշ չափով շփվող/աշխույժ
		\item շփվող/աշխույժ
	\end{enumerate}
	\item Մարդկանց հետ փոխհամաձայնության գալիս ո՞ր տարբերակն է ավելի բնորոշ Ձեզ։
	\begin{enumerate}
		\item մրցակցող/հակաճառող
		\item որոշ չափով մրցակցող/հակաճառող
		\item չեզոք
		\item որոշ չափով զիջող/ապրումակցող
		\item զիջող/ապրումակցող
	\end{enumerate}
	\item Ինչպե՞ս կբնութագրեք Ձեր վարքը նյարդայնացնող իրավիճակներում։
	\begin{enumerate}
		\item հավասարակշռված/հանգիստ
		\item որոշ չափով հավասարակշռված/հանգիստ
		\item չեզոք
		\item որոշ չափով զգացմունքային/բռնկվող
		\item զգացմունքային/բռնկվող 
	\end{enumerate}
}
\end{itemize}





\newpage

\section{Das armenische Alphabet}
\label{appendix:alphabeth_am}

Das moderne armenische Alphabet besteht aus 39 Grapheme. Jedes Phonem im modernen Ostarmenisch kann mindestens durch ein Graphem dargestellt werden. Dahingegen befinden sich manche Buchstaben im Alphabet, die einem Phonem entsprechen. \\    

{
\renewcommand{\arraystretch}{1.2}
\centering
\begin{tabular}{@{}|l|l|l|}
	
\toprule
	
\textbf{Gross} & \textbf{Klein} & \textbf{IPA} \\
	
\midrule
		
{\artm A} &	{\artm a} & \textipa{A}\\\hline
{\artm  B} & {\artm b} & \textipa{b}\\\hline
{\artm G} & {\artm g} & \textipa{g}\\\hline
{\artm D} & {\artm d} & \textipa{d}\\\hline
{\artm E} & {\artm e} & \textipa{E} (\textipa{jE})\\\hline
{\artm Z} & {\artm z} & \textipa{z}\\\hline
{\artm E'} & {\artm e'} & \textipa{E}\\\hline
{\artm U'} & {\artm u'} & \textipa{@}\\\hline
{\artm T'} & {\artm t'} & \textipa{t\super{h}}\\\hline
{\artm Z'} & {\artm g'} & \textipa{Z}\\\hline
{\artm I} & {\artm I} & \textipa{i}\\\hline
{\artm L} & {\artm l} & \textipa{l}\\\hline
{\artm X} & {\artm x} & \textipa{X}\\\hline
{\artm C'} & {\artm c'} & \texttslig\\\hline
{\artm K} & {\artm k} & \textipa{k}\\\hline
{\artm H} & {\artm h} & \textipa{h}\\\hline
{\artm DZ} & {\artm dz} & \textdzlig \\\hline
{\artm GH} & {\artm gh} & \textipa{K}\\\hline
{\artm J'} & {\artm j'} & \textteshlig \\\hline
{\artm M} & {\artm m} & \textipa{m}\\\hline
{\artm Y} & {\artm y} & \textipa{j}\\\hline
{\artm N} & {\artm n} & \textipa{n}\\\hline
{\artm SH} & {\artm sh} & \textipa{S}\\\hline
{\artm O} & {\artm o} & \textipa{O} (\textipa{vO})\\\hline
{\artm CH} & {\artm ch} & \textteshlig{\super{h}} \\\hline
{\artm P} & {\artm p} & \textipa{p}\\\hline
{\artm J} & {\artm j} & \textdyoghlig\\\hline
{\artm R} & {\artm r} & \textipa{r}\\\hline
{\artm S} & {\artm s} & \textipa{s}\\\hline
{\artm V} & {\artm v} & \textipa{v}\\\hline
{\artm T} & {\artm t} & \textipa{t}\\\hline
{\artm R} & {\artm r} & \textfishhookr \\\hline
{\artm C} & {\artm c} & \texttslig{\super{h}}\\\hline
{\artm Ow} & {\artm ow} & \textipa{u}\\\hline
{\artm PH} & {\artm ph} & \textipa{p\super{h}}\\\hline
{\artm Q} & {\artm q} & \textipa{k\super{h}}\\\hline
 & {\artm ev} & \textipa{Ev} (\textipa{jEv})\\\hline
{\artm O'} & {\artm o'} & \textipa{O}\\\hline
{\artm F} & {\artm f} & \textipa{f}\\\hline
	
\bottomrule

\end{tabular}

}

\noindent Manche Grapheme haben mehr als eine Ausspracheoption. Die Aussprache der Wortinitialen Buchstaben entspricht der in Klammern. 

Im Armenischen gibt es kein Großbuchstabe für {\artm ev} [\textipa{Ev}] , [\textipa{jE}], stattdessen werden die Buchstaben {\artm E} [\textipa{E}] , [\textipa{jEv}] und {\artm v} [\textipa{v}]  benutzt. Beispiel:

\end{document}
